\documentclass[12pt]{article}

% --- Packages ---
\usepackage[utf8]{inputenc}
\usepackage[english]{babel}
\usepackage{geometry}
\geometry{a4paper, margin=1in, headheight=15pt}
\usepackage{titlesec}
\usepackage{array}
\usepackage{longtable}
\usepackage{fancyhdr}
\usepackage{hyperref}
\usepackage{enumitem}

% --- Header and Footer ---
\pagestyle{fancy}
\fancyhf{}
\lhead{}
\rhead{Software Requirements Specifications}
\cfoot{Page \thepage}
% --- Custom formatting ---
\titleformat{\section}{\normalfont\Large\bfseries}{\thesection.}{0.5em}{}
\titleformat{\subsection}{\normalfont\large\bfseries}{\thesubsection.}{0.5em}{}

\begin{document}

% --- Title Page ---
\begin{titlepage}
    \centering
    \vspace*{0.5in}
    {\large \textbf{National University of Computer and Emerging Sciences}} \\
    \vspace{1in}
    {\Huge \textbf{Software Requirement Specifications}} \\
    \vspace{0.5in}
    {\LARGE \textbf{Autonomous Multimodal AI Proxy Agent \\ for Automated Class Participation \& Attendance}} \\
    \vspace{1in}
    \textbf{Project Team Members:} \\
    \vspace{0.2in}
    \begin{tabular}{ll}
        24K-1033  \\
        24K-0851 \\
        24K-0956 \\
    \end{tabular} \\
    \vspace{1in}
    \textbf{Supervision:} \\
    Supervisor: Ms. Maarukh \\
    \vfill
    Submission Date: \today
\end{titlepage}

\newpage

% --- Document History ---
\section*{Document History}
\begin{tabular}{|p{2cm}|p{4cm}|p{3cm}|p{5cm}|}
\hline
\textbf{Version} & \textbf{Name} & \textbf{Date} & \textbf{Description} \\ \hline
01.00 & 24K-1033, 24K-0851, 24K-0956 & \today & Initial SRS draft matching Interaction Diagram v1.0 \\ \hline
\end{tabular}

\newpage
\tableofcontents
\newpage

% --- Section 1: Introduction ---
\section{Introduction}

\subsection{Purpose of Document}
This document specifies the requirements for the Autonomous AI Proxy Agent. The system leverages a multimodal architecture to monitor virtual classrooms via browser automation, interpret auditory and textual cues using LLMs, and autonomously respond to instructors. By providing a session link and time, the agent performs a full "proxy" of the student's presence, ensuring attendance is marked and queries are addressed without manual intervention.

\subsection{Intended Audience}
\begin{itemize}
    \item \textbf{Project Supervisor (Ms. Maarukh):} To verify the technical logic of the multimodal pipeline and exception handling modules.
    \item \textbf{Development Team:} To implement the orchestration between the \texttt{BotController}, \texttt{BrowserAutomation}, and AI APIs (Gemini/Groq).
    \item \textbf{System Architects:} To ensure the stability of the \texttt{NotificationManager} and \texttt{ExceptionHandler} during long-duration sessions.
\end{itemize}

\subsection{Abbreviations}
\begin{itemize}
    \item \textbf{LLM}: Large Language Model (Gemini/Groq)
    \item \textbf{STT}: Speech-to-Text (WhisperAI)
    \item \textbf{BA}: Browser Automation (Chromium-based)
    \item \textbf{VAD}: Voice Activity Detection
    \item \textbf{JSON}: JavaScript Object Notation (for \texttt{schedule.json} configuration)
\end{itemize}

% --- Section 2: Overall System Description ---
\section{Overall System Description}

\subsection{Project Background}
As online education becomes standardized, students often face scheduling conflicts. This project addresses the need for a "Digital Proxy"—an agent capable of not just joining a meeting, but actively listening for keywords (like the student's name), analyzing chat sentiment, and generating contextual responses that mimic human interaction.

\subsection{Project Scope}
The system architecture, as defined in the interaction diagram, encompasses:
\begin{itemize}
    \item \textbf{Automated Orchestration}: A \texttt{Scheduler} that polls \texttt{schedule.json} to trigger \texttt{BrowserAutomation} at specific class times.
    \item \textbf{Multimodal Perception}: A \texttt{VoiceListener} using WhisperAI for real-time transcription and a \texttt{ChatAnalyzer} for monitoring textual interactions.
    \item \textbf{Cognitive Reasoning}: An \texttt{AIResponder} that interfaces with Gemini and Groq APIs to generate context-aware replies when triggers are met.
    \item \textbf{Attendance Management}: An \texttt{AttendanceManager} that evaluates confidence levels (90-95\% accuracy) before sending a "Present" confirmation.
    \item \textbf{Remote Telemetry}: A \texttt{NotificationManager} that sends real-time alerts to the user via \texttt{TelegramBotAPI} regarding attendance status or class entry/exit.
\end{itemize}

\subsection{User Classes and Characteristics}
The system is designed for students who require a reliable, hands-free presence in digital environments. Users interact primarily through configuration files (\texttt{schedule.json}) and receive feedback via Telegram.

% --- Section 3: External Interface Requirements ---
\section{External Interface Requirements}

\subsection{User Interfaces}
\begin{itemize}
    \item \textbf{CLI / Startup Script}: The user initiates the bot via \texttt{run.bat}, which handles setup verification and credential loading.
    \item \textbf{Telegram Interface}: The primary mobile interface for receiving "Attendance Marked" notifications or "Class Missed" alerts.
\end{itemize}

\subsection{Software Interfaces}
\begin{itemize}
    \item \textbf{Browser Automation Layer}: Uses Chromium to launch meeting URLs, toggle microphones, and manage browser captions.
    \item \textbf{AI Reasoning APIs}: Dual-integration with \textbf{GeminiAPI} (primary) and \textbf{GroqAPI} (fallback) for prompt-based response generation.
    \item \textbf{Audio Processing}: Integration with \textbf{WhisperAI} for transcribing live classroom audio streams into text for LLM analysis.
\end{itemize}

\subsection{Hardware Interfaces}
\begin{itemize}
    \item \textbf{Virtual Audio Pipeline}: Requires a virtual audio driver to route system/browser output into the \texttt{VoiceListener} module.
    \item \textbf{Network}: Persistent high-bandwidth connection to maintain the browser session and simultaneous API calls to LLM providers.
\end{itemize}

\newpage

\newpage

% --- Section 4: System Modeling & Use Cases ---
\section{System Modeling and Use Cases}

\subsection{Use Case Diagram Overview}
The system's interactions are categorized into 7 primary functional use cases derived from the system architecture. These involve interactions between the \textbf{Student (User)}, the \textbf{Browser Automation Layer}, the \textbf{Multimodal Perception Pipeline (WhisperAI/Chat)}, and external actors like the \textbf{LLM Providers (Gemini/Groq)} and \textbf{Telegram API}.

\vspace{1cm}
\begin{figure}[h!]
    \centering
    \framebox[\textwidth]{\vbox to 9cm{\vfill \centering
        \textbf{UML Use Case Diagram: Autonomous AI Proxy Agent} \\
        \vspace{0.5cm}
        \begin{itemize}
            \item \textbf{UC-01:} Automated Session Scheduling \& Launch
            \item \textbf{UC-02:} Multimodal Classroom Monitoring (Audio/Chat)
            \item \textbf{UC-03:} Attendance Verification \& Response Logic
            \item \textbf{UC-04:} Contextual Query Processing (LLM Reasoning)
            \item \textbf{UC-05:} API Failover \& Redundancy Management
            \item \textbf{UC-06:} Remote Telemetry \& Telegram Notification
            \item \textbf{UC-07:} Session Exit \& Resource Cleanup
        \end{itemize}
    \vfill}}
    \caption{Use Case Diagram showing the flow from schedule initialization to autonomous interaction and notification.}
    \label{fig:use_case}
\end{figure}
\vspace{1cm}

\subsection{Detailed Use Case Tables}

% --- Use Case 1 ---
\subsubsection{UC-01: Automated Session Scheduling \& Launch}
\begin{longtable}{|p{4cm}|p{10cm}|}
\hline
\textbf{Use Case ID} & UC-01 \\ \hline
\textbf{Primary Actor} & Student (User) \\ \hline
\textbf{Pre-conditions} & \texttt{schedule.json} is configured; \texttt{run.bat} executed. \\ \hline
\textbf{Main Flow} & 1. Scheduler polls the JSON file for class time. 2. BotController triggers BrowserAutomation. 3. System launches Chromium and joins meeting via URL. \\ \hline
\end{longtable}

% --- Use Case 2 ---
\subsubsection{UC-02: Multimodal Classroom Monitoring}
\begin{longtable}{|p{4cm}|p{10cm}|}
\hline
\textbf{Use Case ID} & UC-02 \\ \hline
\textbf{Primary Actor} & VoiceListener / ChatAnalyzer \\ \hline
\textbf{Secondary Actor} & WhisperAI \\ \hline
\textbf{Main Flow} & 1. VoiceListener captures raw audio stream. 2. WhisperAI transcribes audio to text. 3. ChatAnalyzer scans meeting chat for name/keywords. \\ \hline
\end{longtable}

% --- Use Case 3 ---
\subsubsection{UC-03: Attendance Verification \& Response}
\begin{longtable}{|p{4cm}|p{10cm}|}
\hline
\textbf{Use Case ID} & UC-03 \\ \hline
\textbf{Primary Actor} & AttendanceManager \\ \hline
\textbf{Main Flow} & 1. System detects "attendance" or "name" trigger. 2. AttendanceManager evaluates trigger accuracy (90-95\% confidence). 3. System types "Present" or triggers vocal confirmation. \\ \hline
\end{longtable}

% --- Use Case 4 ---
\subsubsection{UC-04: Contextual Query Processing}
\begin{longtable}{|p{4cm}|p{10cm}|}
\hline
\textbf{Use Case ID} & UC-04 \\ \hline
\textbf{Primary Actor} & AIResponder \\ \hline
\textbf{Secondary Actor} & Gemini / Groq API \\ \hline
\textbf{Main Flow} & 1. System identifies a question directed at the student. 2. Text/Context sent to LLM. 3. LLM generates a human-like response for the meeting. \\ \hline
\end{longtable}

% --- Use Case 5 ---
\subsubsection{UC-05: API Failover Management}
\begin{longtable}{|p{4cm}|p{10cm}|}
\hline
\textbf{Use Case ID} & UC-05 \\ \hline
\textbf{Primary Actor} & ExceptionHandler \\ \hline
\textbf{Main Flow} & 1. Primary API (Gemini) fails or returns error. 2. ExceptionHandler catches error and logs to file. 3. System automatically switches to GroqAPI to maintain session. \\ \hline
\end{longtable}

% --- Use Case 6 ---
\subsubsection{UC-06: Remote Telemetry \& Notification}
\begin{longtable}{|p{4cm}|p{10cm}|}
\hline
\textbf{Use Case ID} & UC-06 \\ \hline
\textbf{Primary Actor} & NotificationManager \\ \hline
\textbf{Secondary Actor} & Telegram API \\ \hline
\textbf{Main Flow} & 1. System confirms attendance marked or class joined. 2. NotificationManager formats a status alert. 3. User receives real-time update on Telegram mobile app. \\ \hline
\end{longtable}

% --- Use Case 7 ---
\subsubsection{UC-07: Session Exit \& Resource Cleanup}
\begin{longtable}{|p{4cm}|p{10cm}|}
\hline
\textbf{Use Case ID} & UC-07 \\ \hline
\textbf{Primary Actor} & BotController \\ \hline
\textbf{Post-conditions} & Browser process killed; network resources released. \\ \hline
\textbf{Main Flow} & 1. Exit condition met (e.g., participants < 5 or class time elapsed). 2. BotController sends leave command. 3. System returns to polling state for next class. \\ \hline
\end{longtable}

\newpage
% --- Section 4: Functional Requirements ---
\section{Functional Requirements}

\subsection{Session Automation and Control}
\begin{itemize}
    \item \textbf{FR-01: Automated Scheduling}: The system shall poll the \texttt{schedule.json} file to identify upcoming sessions and trigger the \texttt{BotController} at the exact timestamp.
    \item \textbf{FR-02: Browser Orchestration}: The system must autonomously launch a Chromium instance, navigate to meeting URLs, and manage UI elements (toggling mic/camera and enabling captions).
    \item \textbf{FR-03: Dynamic Exit Logic}: The system shall monitor the participant count and automatically terminate the session if the count drops below 5 or if the scheduled duration plus a 5-minute buffer is reached.
\end{itemize}

\subsection{Multimodal Perception and Response}
\begin{itemize}
    \item \textbf{FR-04: Real-time Audio Transcription}: The system shall utilize WhisperAI to transcribe live audio streams into text with a focus on detecting "name" or "attendance" keywords.
    \item \textbf{FR-05: Collaborative Chat Analysis}: The system must monitor meeting chat boxes for specific patterns and high-frequency "Present" bursts to trigger attendance responses.
    \item \textbf{FR-06: Dual-LLM Reasoning}: The \texttt{AIResponder} shall generate context-aware answers to teacher queries by querying GeminiAPI (primary) or GroqAPI (fallback).
\end{itemize}

\subsection{Attendance and Notification}
\begin{itemize}
    \item \textbf{FR-07: Attendance Verification}: The system must evaluate triggers with an accuracy threshold of 90-95\% before executing an "Attendance Marked" action.
    \item \textbf{FR-08: Remote Status Updates}: The system shall send real-time event logs (e.g., "Class Joined," "Attendance Marked," "Class Missed") to the user via the \texttt{TelegramBotAPI}.
\end{itemize}

% --- Section 5: Non-Functional Requirements ---
\section{Non-Functional Requirements}

\subsection{Performance and Latency}
\begin{itemize}
    \item \textbf{Processing Speed}: The conversion from raw audio stream to WhisperAI transcription must occur within a sub-2.0 second window to ensure response relevance.
    \item \textbf{API Latency}: The failover mechanism between Gemini and Groq must complete within 1.5 seconds to prevent awkward silences in the meeting.
    \item \textbf{Resource Footprint}: The browser automation must run efficiently on systems with 4GB to 8GB of RAM without causing audio stuttering.
\end{itemize}

\subsection{Reliability and Resilience}
\begin{itemize}
    \item \textbf{Redundancy}: If the primary LLM API fails, the \texttt{ExceptionHandler} must successfully resume the query using the secondary provider without session termination.
    \item \textbf{Network Recovery}: The system shall attempt to rejoin the meeting automatically if a temporary network failure is detected by the \texttt{Scheduler}.
    \item \textbf{Logging}: All errors, API failures, and transaction timestamps must be logged to \texttt{log.txt} for post-session auditing.
\end{itemize}

\subsection{Security and Data Integrity}
\begin{itemize}
    \item \textbf{Credential Isolation}: Meeting links and API keys must be loaded from local environment configurations to prevent data leakage.
    \item \textbf{Zero-Trust Communication}: Data sent to the \texttt{AIResponder} must be sanitized to prevent prompt injection from meeting chat participants.
\end{itemize}

\subsection{Resource Management}
To optimize system health, the "End all connections" protocol must be strictly enforced. Upon meeting exit, the system must kill all Chromium sub-processes and close API sockets to prevent memory leaks and unnecessary billing from LLM providers.

\end{document}
